% chapter2.tex
\chapter{Research and Literature Review}

The development of a system that converts hand-drawn electronic circuit diagrams into machine-readable simulation formats draws on electronic design automation (EDA), computer vision, object detection, and document understanding. This chapter surveys the relevant literature, compares available datasets and detection architectures, and honestly assesses the gaps this project addresses.

\section{Circuit Digitization: Motivation and Approaches}

Electronic circuit diagrams are the primary medium between a designer's intent and the simulation tools that validate it. Despite sophisticated EDA tools such as LTspice, KiCad, and Cadence Virtuoso, hand-drawn sketches remain the dominant medium for rapid conceptualization because they impose no syntactic overhead. The problem is that these sketches exist in an analogue domain inaccessible to simulators, forcing engineers to manually re-enter topologies—a time-consuming, error-prone process.

\begin{figure}[H]
    \centering
    \includegraphics[width=0.95\textwidth]{images/schemdraw_example.png}
    \caption{An example of programmatic circuit generation using the SchemDraw library, demonstrating the syntax-heavy nature of manual digitization.}
    \label{fig:schemdraw}
\end{figure}



Research in diagram recognition bifurcates into \emph{online} and \emph{offline} approaches. Online systems process pen-stroke sequences as they are generated, exploiting rich temporal information but requiring specialised hardware. Offline systems treat the diagram as a raster image and must recover all structural information from pixel data alone—the more practically relevant paradigm for retrofitting existing paper records. The system described in this report operates entirely offline: the user uploads a photograph of a hand-drawn circuit, and the pipeline identifies components and generates a partial netlist.

\section{Dataset Comparison}

\subsection{Kaggle Handwritten Circuit Dataset}

The primary training resource for this project is a Kaggle dataset of annotated hand-drawn circuit images spanning 9 classes: passive components (resistors, capacitors, inductors), active devices (MOSFETs, BJTs, op-amps), and structural elements (junctions, wire segments, text bounding boxes). The images exhibit natural variation in pen weight, paper texture, and lighting. Several limitations directly constrain the project's scope: the class distribution is significantly imbalanced, illumination variance complicates feature extraction, and—most critically—the dataset provides bounding-box annotations for text regions but does not transcribe their content. Consequently, the trained model can localise a label such as ``10k$\Omega$'' but cannot extract its value, scoping the current implementation to \emph{topological} component detection only.

\begin{figure}[H]
    \centering
    \includegraphics[width=0.9\textwidth]{images/sample_from_dataset.jpg}
    \caption{Sample images from the Kaggle dataset showcasing varying sketching styles, illumination, and component orientations.}
    \label{fig:dataset_samples}
\end{figure}

\subsection{Thoma et al.\ and Digitize-HCD Datasets}

Thoma et al.~\cite{thoma2021} released one of the earliest public datasets targeting offline handwritten circuit diagram recognition: 1,152 images from 144 unique circuits drawn by 12 drafters, yielding 48,563 annotations across 45 classes under varied capture conditions including whiteboards and aluminium foil. Critically, the dataset introduces \emph{auxiliary annotation classes}—junctions, crossovers, and terminal endpoints—necessary to distinguish a T-junction (three connected wires) from a crossover (two visually overlapping but electrically isolated wires). The Kaggle dataset lacks systematic crossover annotation, which is a primary reason connectivity extraction remains unreliable in the current implementation.

Ahmed et al.~\cite{ahmed2025} present Digitize-HCD, the most comprehensive public resource currently available: 1,277 images from 176 volunteers, annotated at three granularity levels—component bounding boxes across 17 symbol classes, transcribed text labels using polygon annotations, and component \emph{port localization} via probability heatmaps encoding precise terminal coordinates. The port heatmaps address the hardest sub-problem in circuit digitization: knowing not just where a component is but where its pins are, and hence which wires connect to which terminal. The Kaggle dataset provides no port annotations, meaning the current implementation cannot determine transistor terminal connectivity without additional heuristics. Digitize-HCD's text sub-dataset further includes transcribed strings for 11,936 labels, enabling OCR training—no equivalent transcription exists in the Kaggle dataset.

\section{Object Detection Architectures}

\subsection{Template Matching, Faster R-CNN, and SSD}

Template matching—sliding a reference image across the input and computing similarity at each position—requires no training data but proves entirely inadequate for hand-drawn circuits. The same resistor drawn by two individuals may differ substantially in size, stroke style, and orientation, causing similarity scores to collapse and generalisation to fail.

Faster R-CNN~\cite{ren2015} decomposes detection into a Region Proposal Network (RPN) that generates candidate bounding boxes followed by per-region classification on shared convolutional features. It achieves high localisation precision; Thoma et al.\ establish a 52\% mAP baseline on their dataset using Faster R-CNN with a ResNet-152 backbone. However, its sequential two-stage pipeline carries substantial computational overhead and complexity that presents a practical barrier for a student-level prototype.

SSD~\cite{liu2016} collapses both stages into a single forward pass by predicting boxes and class probabilities simultaneously at multiple feature map scales using fixed anchor boxes. While faster than Faster R-CNN, SSD exhibits a known weakness with very small or densely packed objects—a significant limitation for circuit diagrams where junction dots and terminal endpoints occupy only a handful of pixels.

\subsection{YOLO: Unified Real-Time Detection}

The YOLO architecture~\cite{redmon2016} reframes detection as a single regression task. The input image is divided into an $S \times S$ grid; each cell predicts $B$ bounding boxes and $C$ conditional class probabilities in a single forward pass, producing an $S \times S \times (B \cdot 5 + C)$ output tensor. YOLO is trained end-to-end by a multi-part loss function that jointly penalises coordinate, confidence, and classification errors. To equalise gradient contributions across object sizes, bounding box dimensions are predicted as square roots of width and height. Non-maximal suppression resolves overlapping predictions at inference time.

The principal advantage is speed: the base model achieves 45 frames per second on GPU hardware, and Fast YOLO exceeds 150 fps while outperforming prior real-time detectors by more than double the mAP~\cite{redmon2016}. YOLO also reasons from global image context, reducing spurious false positives in structured images. Its known limitations—grid-cell density constraints making dense small-object detection difficult, and higher localisation error than two-stage detectors—are relevant to tightly packed components and small structural elements in circuit diagrams.

\subsection{Justification for YOLOv8}

YOLOv8 was selected following a comparative assessment of the available architectures against the project's practical constraints. Its most significant advancement over earlier YOLO versions is an \emph{anchor-free} detection head that predicts boxes as direct regression targets from learned centre-point predictions, eliminating the dataset-specific anchor tuning required by anchor-based detectors—particularly advantageous when the same component class appears at dramatically different aspect ratios across drafters. YOLOv8 is implemented within the Ultralytics/PyTorch framework, provides pre-trained COCO weights for transfer learning, exports inference-ready \texttt{.pt} files for fully offline execution, and exposes a clean Python API that allowed engineering effort to focus on higher-level connectivity and netlist challenges rather than low-level model construction.

\section{The Circuit Understanding Problem}

\subsection{Connectivity Extraction}

A common misconception is that component detection is the primary challenge. In reality, \emph{connectivity extraction}—determining which components are electrically connected and through which terminals—is the deeper and harder problem. Without it, a list of detected components is merely an inventory and cannot be converted into a valid SPICE netlist. The distinction is concrete: a detector that correctly identifies both resistors in a voltage divider captures the component inventory, but the netlist additionally requires knowing which terminal connects to the supply node, which to the shared mid-node, and which to ground. Connectivity extraction is therefore a graph inference problem, not a classification problem.

The standard approach applies thinning or skeletonisation to a binarised image, followed by graph analysis of the resulting line skeleton to identify junctions. This faces several practical difficulties. Junction detection requires pixel-level precision; bounding-box detectors may localise junction dots too imprecisely for correct graph analysis. The \emph{node merging problem} arises at complex intersections where three or more wires converge, requiring the algorithm to correctly count distinct nodes in noisy hand-drawn images. Most critically, distinguishing genuine junctions from wire crossovers is non-trivial without explicit crossover annotation—an absence in the Kaggle dataset that leaves the current implementation without a principled mechanism for this distinction.

Wire interpretation itself is non-trivial: hand-drawn wires may be slightly curved, discontinuous at pen lifts, or confused with component symbol line segments. Skeletonisation produces spurious branch points in noisy ink regions, creating phantom junctions. The current implementation uses basic morphological heuristics rather than a learned wire segmentation model, causing quality to vary significantly with lighting, pen weight, and paper texture.

\subsection{NGSPICE Netlist Generation}

NGSPICE accepts circuit descriptions as plain-text netlists following the convention \texttt{[Identifier] [Node\_1] [Node\_2] [Value]}, where the identifier's leading letter encodes component type (R, C, L, V, Q, M, etc.). Generating a valid netlist requires a three-stage pipeline: (i) localising each component's spatial extent from bounding boxes; (ii) assigning shared node labels to terminals connected by the same wire; and (iii) transcribing each component into netlist syntax with its class, node assignments, and value.

The current implementation completes stage~(i) reliably for well-drawn images. Stage~(ii) is only partially implemented: a proximity-based terminal merging algorithm produces plausible results for simple circuits but fails at complex junctions and crossovers, leading to node merging errors—the most common failure mode observed in testing. Stage~(iii) is consequently impaired by stage~(ii) errors, and because no OCR is implemented, all component values appear as placeholder tokens. The output is therefore a \emph{partial netlist}: it cannot be fed directly to NGSPICE without manual correction, but serves as a structured intermediate representation that preserves the component inventory and a coarse topology approximation. For simple series or parallel circuits it may require only minor edits; for circuits involving active devices or dense junctions, significant human intervention is required. This characterises the current development stage honestly rather than as a failure.

\section{Research Gaps and Summary}

The literature review reveals several persistent gaps. First, fully integrated end-to-end systems that accept a raw photograph and output a directly simulable netlist remain largely absent from public literature; most published systems demonstrate competence in one or two pipeline stages and defer the rest as future work. Second, even in comprehensive datasets like Digitize-HCD, integration of OCR results into downstream connectivity and netlist stages has not been systematically demonstrated—associating a recognised value string with the correct component requires spatial reasoning that is non-trivial across varied drawing styles. Third, connectivity extraction—determining circuit graph structure from wire geometry—remains broadly acknowledged as the hardest unsolved sub-problem; the port localization heatmaps of Ahmed et al.\ represent a promising step, but robust wire segmentation in noisy images remains an active challenge. Fourth, most published systems originate in well-resourced research settings; demonstrating meaningful progress within undergraduate project constraints is itself a reproducibility contribution.

In summary, the Kaggle dataset provides a serviceable training resource for component detection but falls short of Thoma et al.\ and Digitize-HCD in junction annotation, port localization, and transcribed text. YOLOv8 is justified by its speed, anchor-free head, and practical implementability, while its limitations for dense small-object scenes are acknowledged. Component detection is reliably operational; OCR is not implemented; connectivity extraction is preliminary and error-prone; and netlist generation produces partial, manually-correctable output. Each of these limitations defines a concrete direction for future work, and the system provides a foundation upon which those extensions can be built.

